\section{From Recurring Failure to a Procedure Bottleneck}
\label{sec:mechanism}
Let endpoint $i$ specify a question $x_i$, a cutoff-valid evidence package $E_i$, a model configuration $M_i$, and a family-specific binary success rule $Y_i\in\{0,1\}$. We intervene only on the reusable procedure exposed to the frozen harness. The three conditions are $N$ (no experimental helper), $G$ (a complexity-matched target-blind helper), and $T$ (the mechanism-specific targeted procedure). All other endpoint inputs and scoring rules are held fixed.

For a set of frozen endpoints $\mathcal I$ within one model--task regime, we report two empirical contrasts,
\begin{equation}
\widehat{\Delta}_{T-N}
= \frac{1}{|\mathcal I|}\sum_{i\in\mathcal I}\left(\bar Y_i(T)-\bar Y_i(N)\right),
\qquad
\widehat{\Delta}_{T-G}
= \frac{1}{|\mathcal I|}\sum_{i\in\mathcal I}\left(\bar Y_i(T)-\bar Y_i(G)\right),
\end{equation}
where $\bar Y_i(a)$ averages repeated executions of the same endpoint under arm $a$. Repeats characterize execution robustness and do not increase the number of independent endpoint units.

\paragraph{Definition 1 (intervention-defined procedure bottleneck).}
A targeted procedure is a \emph{binding repair} for a frozen model--task regime only when it improves the registered family outcome over both the original agent and matched generic assistance, i.e., when the endpoint evidence supports $\widehat{\Delta}_{T-N}>0$ and $\widehat{\Delta}_{T-G}>0$. A positive targeted-minus-generic contrast with $\widehat{\Delta}_{T-N}\leq0$ is explicitly classified as \emph{no repair}. This definition identifies the value of inserting the operation under the frozen intervention; it does not imply causal exclusivity or rule out another implementation of the same operation.

\begin{figure}[t]
\centering
\includegraphics[width=\linewidth]{figures/fig1_temporal_skill_causal.pdf}
\caption{Three-arm procedure-bottleneck test. The same frozen endpoint is evaluated as the original no-skill agent (N), with a complexity-matched target-blind helper (G), and with the targeted reusable procedure (T). A binding repair requires T to beat both N and G; unchanged-code held-out replay separately tests reuse.}
\label{fig:causal}
\end{figure}

The design targets three failure modes that recur as release content changes.

\subsection{Cutoff validity}
The cutoff procedure receives a query date and candidate evidence timestamps and returns only evidence admissible by that date. It addresses future leakage: a later release may be highly relevant yet invalid for an earlier information set. The matched generic helper can validate evidence-package structure but is forbidden from comparing release dates with the cutoff.

\subsection{Release alignment}
The alignment procedure maps a source release to the reporting period represented by its content and verifies the selected time-series row against that period. It addresses cases in which publication time and reporting time differ. The matched generic helper can inspect document structure but cannot perform the date-to-period mapping.

\subsection{Exogenous grounding}
The grounding procedure links a numerical change to first-party documentary evidence available before the cutoff. It operates over the full frozen candidate-span region rather than a gold-only pool and returns evidence tied to the requested attribution slots. Its generic helper exposes document structure without the mechanism-specific contextual event selector.

No targeted procedure contains scenario-specific answers, gold labels, or endpoint-specific fitted values. Source hashes are frozen before execution, and unchanged implementations are reused on held-out periods. This gives three direct tests of the bottleneck interpretation: a two-arm false positive must be rejected by N; gains should disappear at no-skill ceiling; and compatible held-out releases should preserve direction without rewriting the procedure.
